% Produktivierung — claude-tunnel bachelor thesis.
\chapter{Produktivierung}
\label{ch:produktivierung}

Der in Kapitel~\ref{ch:implementierung} beschriebene Prototyp weist die
Machbarkeit des providerblinden Tunnels nach, ist als Testbett aber noch kein
betreibbarer Dienst: der Zustand liegt ausschliesslich im Speicher, es gibt
keine konventionelle Nutzeridentitaet, und die Zertifikate sind fest
verdrahtet. Dieses Kapitel beschreibt die Ueberfuehrung in ein
produktionsreifes System entlang sechs Bausteinen -- Persistenz
(Abschnitt~\ref{sec:prod-persistenz}), Identitaet
(\ref{sec:prod-identitaet}), PKI und Transportsicherheit
(\ref{sec:prod-pki}), Auslieferung (\ref{sec:prod-deploy}), Haertung
(\ref{sec:prod-haertung}) und Bezahlung (\ref{sec:prod-payment}). Die
kryptografische Ende-zu-Ende-Vertraulichkeit der Nutzlast (Noise,
\cite{perrin2018noise}) bleibt in allen Bausteinen erhalten: die
Produktivierung fuegt Identitaet und Abrechnung hinzu, ohne dem Betreiber
Zugriff auf den Klartext zu geben.

\section{Durabler Zustand}
\label{sec:prod-persistenz}

Enrollment, Tunnel-Registry und Guthaben-Ledger wurden von
In-Memory-Strukturen auf eine eingebettete SQLite-Datenbank umgestellt. Jeder
Speicher kapselt seine Tabellen hinter einer synchron gehaltenen Verbindung;
mehrstufige Operationen wie die Zahlungsbestaetigung laufen in einer
Transaktion, sodass ein Absturz weder ein Guthaben doppelt gutschreibt noch
einen Halbzustand hinterlaesst. Die Durabilitaet ist durch Neustart-Tests
abgesichert: eine zweite Dienstinstanz auf derselben Datei sieht denselben
Zustand, ein bereits eingeloester Join-Token bleibt verbraucht.

\section{Identitaet und Authentifizierung}
\label{sec:prod-identitaet}

Statt pseudonymer Konten treten konventionelle Konten ueber einen
OIDC-Anbieter (Keycloak). Der Kontrollknoten verifiziert Bearer-Token per
RS256 gegen den oeffentlichen Realm-Schluessel und prueft Aussteller und
Ablauf; das Konto wird aus dem Token-Subjekt abgeleitet, sodass ein Aufrufer
stets nur auf sein eigenes Konto wirkt. Bewusst wird damit die frühere
Pseudonymitaets-Aussage aufgegeben -- der Betreiber kennt die Identitaet zur
Abrechnung. Die inhaltliche Vertraulichkeit bleibt hingegen bestehen: der
Datenpfad ist unveraendert Ende-zu-Ende verschluesselt.

\section{PKI und Transportsicherheit}
\label{sec:prod-pki}

Der Edge betreibt eine interne Zertifizierungsstelle. Clients vertrauen der
CA-Wurzel statt einem angehefteten Blatt-Zertifikat; dadurch kann der Edge
sein Zertifikat rotieren, ohne dass ein Client geaendert werden muss. Auf der
Transportebene gilt \emph{TLS ueberall}: QUIC mit TLS~1.3
(\cite{rfc9001,rfc8446}) und ein TLS-ueber-TCP-Rueckfallpfad zum Edge, HTTPS
am Ingress vor dem Kontrollknoten. Der einzige im Klartext sprechende
Bestandteil -- die HTTP-Schnittstelle des Kontrollknotens -- verlaesst den
Cluster nie.

\section{Auslieferung}
\label{sec:prod-deploy}

Der Dienst ist auf zwei Wegen auslieferbar. Fuer den Eigenbetrieb genuegt ein
gehaertetes Compose-Buendel mit persistentem Datenvolumen, Neustart-Richtlinie
und einer Bereitschaftspruefung auf \texttt{/readyz}. Fuer den gehosteten
Betrieb existieren Kubernetes-Manifeste (kustomize) mit einem PVC-gestuetzten
Kontrollknoten, Liveness-/Readiness-Sonden, gehaerteten Sicherheitskontexten
und einem TLS-terminierenden Ingress. Beide Wege nutzen dieselben Binaries und
dasselbe Protokoll.

\section{Haertung}
\label{sec:prod-haertung}

Neben dem Proof-of-Work-Gatter (\cite{back2002hashcash}) begrenzt ein
Ratenlimit je Konto die Token-Ausgabe, sodass ein einzelnes Konto den
Kontrollknoten nicht erschoepfen kann. Die Abhaengigkeiten werden gegen die
RustSec-Advisory-Datenbank geprueft und ueber eine eingecheckte
\texttt{Cargo.lock} exakt gepinnt; ein Wächter blockiert versehentlich
eingecheckte Geheimnisse. Ein aktualisiertes Bedrohungsmodell haelt fest, was
der Dienst leistet und was nicht -- insbesondere macht er keine
Anonymitaetsaussage.

\section{Bezahlung}
\label{sec:prod-payment}

Die Guthaben-Gutschrift wird nur noch durch ein signiertes Ereignis des
Zahlungsanbieters ausgeloest, nicht durch einen vom Client aufrufbaren
Endpunkt. Der Kontrollknoten verifiziert die HMAC-SHA256-Signatur des Webhooks
gegen ein geteiltes Geheimnis und prueft die Frische des Zeitstempels gegen
Wiedereinspielung; ohne konfiguriertes Geheimnis wird jeder Webhook abgelehnt
(fail-safe). Die Zustellung ist idempotent, sodass ein wiederholtes Ereignis
nicht doppelt gutschreibt.

\section{Zusammenfassung}
\label{sec:prod-fazit}

Mit durablem Zustand, echter Identitaet, rotierbarer PKI, reproduzierbarer
Auslieferung, Haertung und signaturgesicherter Bezahlung ist aus dem Testbett
ein betreibbarer Dienst geworden -- ohne die zentrale Eigenschaft aufzugeben,
dass der Betreiber die Nutzlast nicht lesen kann.
